\section*{Introduction}

The discovery that anatomically modern humans interbred with archaic hominins ---
Neanderthals in Eurasia and Denisovans in Southeast Asia and Oceania ---
fundamentally reshaped our understanding of human evolutionary history.
Initial estimates from whole-genome sequencing of Vindija Cave Neanderthal
and Altai Denisovan genomes established that $\sim$1--4\% of the genome of
non-African modern humans is derived from archaic
introgression~\citep{green2010draft, reich2010genetic, prufer2014complete}.
Subsequent work demonstrated that these archaic-derived segments are not
evolutionary relics but functionally active contributors to modern human
phenotypic diversity and disease risk: archaic haplotypes have been identified
in immune genes, metabolic pathways, and neurological
circuits~\citep{sankararaman2014genomic, vernot2014resurrecting,
dannemann2017contribution, browning2018ancestry}.
The problem of reliably detecting, mapping, and attributing these introgressed
segments at per-haplotype resolution across large population cohorts therefore
sits at the intersection of evolutionary genomics, population genetics methodology,
and medical genetics.

Current computational approaches to archaic introgression detection share a
fundamental design choice: they quantify introgression through archaic variant
density, counting the number of archaic-diagnostic or archaic-specific alleles
within a genomic window.
HMMix~\citep{skov2018detecting}, DAIseg~\citep{planche2024daiseg}, and related tools
model this density via Poisson or mixture emissions, using the excess of
archaic-matching variants over a background frequency to flag introgressed regions.
IBDmix~\citep{chen2020identifying} extends this principle to identity-by-descent
comparisons at individual sites.
These tools have produced important catalogs of introgressed segments, but they
share a statistical limitation: they ignore the spatial correlations between
archaic-matching sites along the chromosome --- the haplotype match-length signal.

Coalescent theory predicts that this ignored signal is substantial.
An introgressed haplotype entered the modern human gene pool $\sim$50,000 years ago
(approximately 1,724 generations), while a non-introgressed haplotype last shared
a common ancestor with an archaic genome $\sim$550,000 years ago
(approximately 19,000 generations).
The expected length of a perfect haplotype match to the archaic reference therefore
differs by a factor of $\sim$11 between introgressed and non-introgressed segments.
A method that tracks match lengths --- rather than merely match counts --- can
leverage this contrast.
This is precisely the insight behind the Li \& Stephens (LS) haplotype copying
model~\citep{li2003modeling}, which has become the foundation of modern local
ancestry inference tools including ChromoPainter~\citep{lawson2012inference},
MOSAIC~\citep{salter2019mosaic}, and SparsePainter~\citep{yang2025sparsepainter}.
In the modern ancestry setting, the LS model outperforms window-based frequency
methods precisely because haplotype match length is informative about population
of origin~\citep{price2009sensitive}.
Some methods have partially exploited haplotype structure for archaic
introgression: S'~PRIME~\citep{browning2018ancestry} uses LD patterns to
identify introgressed fragments, and ArchIE~\citep{durvasula2019statistical}
uses haplotype features in a logistic regression framework.
However, neither adopts the Li \& Stephens copying model with match-length
emissions as the core inference mechanism, leaving the full potential of
coalescent match-length statistics unexploited for archaic ancestry detection.

Here we introduce \textbf{ArchaicPainter} (\AP{}), a three-state Hidden Markov
Model for archaic introgression detection that adapts the tract-length
transition structure of the Li \& Stephens haplotype copying model to
archaic ancestry inference.
Unlike the LS model, whose hidden states index individual reference haplotypes,
\AP{} defines three hidden states --- anatomically modern human (AMH),
Neanderthal (NEA), and Denisovan (DEN) --- and computes emission probabilities
from population allele frequencies (AMH state) and coalescent match statistics
(archaic states), with transition probabilities parameterized by recombination
distance and expected introgressed segment lengths from published admixture
time estimates.
Crucially, all core HMM parameters in \AP{} (admixture times, split times,
introgression fractions) have explicit biological interpretations and are derived
from published coalescent estimates rather than fitted to training data;
post-processing thresholds (gap-closing and minimum segment length) are calibrated
by a one-time grid search on simulation and held fixed across all real-data analyses.
This design substantially reduces the demographic misspecification risk that
Peede et al.~\citeyear{peede2026review}
identified as a central limitation of simulation-trained machine learning methods
for this problem.

We validate \AP{} in three steps.
On 50 independent simulation replicates with known ground truth
(5~Mb per replicate, msprime~\citep{kelleher2016efficient}),
\AP{} achieves segment-level F1 $= 0.480 \pm 0.095$ (AUPRC $= 0.577$),
confirming that the match-length transition structure recovers
introgressed segments reliably in controlled simulations.
Ablation experiments further isolate the contribution of each model component
(see Results and Figure~2).
Applied to chromosome~21 of the 1000~Genomes Project Phase~3
GRCh38 release for European (CEU) and East Asian (CHB) populations,
\AP{} detects $1.10\% \pm 0.60\%$ Neanderthal ancestry per haplotype
in CEU and $0.98\% \pm 0.52\%$ in CHB, consistent with published
genome-wide estimates of 1--2\%, at a median segment resolution of $\sim$28--57~kb.
As an illustrative application, intersecting \AP{}-detected chr21
introgressed regions with gene annotations identifies 7 of 16 previously
published adaptive introgression candidates and a complete interferon
receptor gene cluster (\textit{IFNAR1}, \textit{IFNAR2},
\textit{IFNGR2}, \textit{IL10RB})~\citep{dannemann2017contribution},
demonstrating that the method recovers biologically meaningful signals
in an unsupervised manner.

The specific contributions of this work are:

\begin{itemize}
  \item A tract-length--parameterized local ancestry HMM for archaic introgression
    that, to our knowledge, is the first method for this problem to use
    coalescent match-length statistics rather than variant density
    as the primary discrimination signal, while explicitly distinguishing
    the roles of admixture time (archaic tract length) and background inter-introgression
    spacing in the transition model.
  \item A principled treatment of reference--donor divergence via a positive-only
    emission mode, with empirical validation of its necessity through simulation ablation.
  \item Simulation benchmarks demonstrating F1 $= 0.480 \pm 0.095$
    (AUPRC $= 0.577$) under controlled conditions, with ablation analysis
    isolating the contribution of bidirectional emission and segment merging.
  \item Proof-of-concept three-state detection simultaneously attributing query
    haplotypes to NEA or DEN ancestry in a two-source admixture setting.
  \item Near-linear computational scaling (empirical exponent 1.21) enabling
    population-scale analysis without specialized hardware.
  \item Open-source software (\AP{}) with fully scripted analysis pipelines
    and a chromosome~21 functional annotation use case demonstrating trait-linked
    archaic variant discovery.
\end{itemize}
